\documentclass[11pt]{article}
\usepackage{amsmath,amssymb,amsthm}
\usepackage{algorithm,algpseudocode}
\usepackage[margin=1in]{geometry}
\usepackage[colorlinks=true,linkcolor=blue!60!black,citecolor=blue!60!black]{hyperref}

\numberwithin{equation}{section}

\newcommand{\tr}{\operatorname{tr}}
\newcommand{\bigo}{\mathcal{O}}
\newcommand{\Ktil}{\widetilde K}
\newcommand{\Amean}{A}

\title{EFGP log-marginal gradient: a step-by-step derivation}
\author{}
\date{}

\begin{document}
\maketitle

\noindent This note derives the EFGP gradient algorithm from scratch. The target
form, which we call the \emph{$\beta$-form}, avoids a $1/\sigma^2$ amplification
of the CG residual that appears in the naive Woodbury expansion. The derivation
is linear: each section uses only what was established in the previous one.

\section{Step 1. The log-marginal and its gradient}

Let $K_\theta\in\mathbb{R}^{N\times N}$ be the kernel Gram matrix, $\sigma^2$ the
observation noise, and $\Ktil := K_\theta + \sigma^2 I$. The Gaussian
log-marginal is
\[
  \log p(y\mid\theta) \;=\; -\tfrac12\,y^\top \Ktil^{-1}y
                          \;-\; \tfrac12\,\log\det\Ktil + \text{const}.
\]
For any scalar hyper $\theta$ (a kernel parameter such as lengthscale or
variance, or the noise $\sigma^2$ itself), the two standard identities
$\partial_\theta\log\det\Ktil = \tr(\Ktil^{-1}\partial_\theta\Ktil)$ and
$\partial_\theta\Ktil^{-1} = -\Ktil^{-1}(\partial_\theta\Ktil)\Ktil^{-1}$ give
\begin{equation}
  \partial_\theta\log p \;=\; \tfrac12\bigl(T_2 - T_1\bigr),\qquad
  T_2 \;=\; y^\top \Ktil^{-1}(\partial_\theta\Ktil)\,\Ktil^{-1}y,\qquad
  T_1 \;=\; \tr\!\bigl(\Ktil^{-1}\partial_\theta\Ktil\bigr).
  \label{eq:gradient}
\end{equation}
Writing $\alpha := \Ktil^{-1}y$,
\begin{equation}
  T_2 \;=\; \alpha^\top(\partial_\theta\Ktil)\,\alpha. \label{eq:T2alpha}
\end{equation}
So we need two primitives: a single solve $\alpha = \Ktil^{-1}y$ (for $T_2$), and
a trace of $\Ktil^{-1}\partial_\theta\Ktil$ (for $T_1$).
For kernel hypers $\partial_\theta\Ktil = \partial_\theta K$; for noise
$\partial_{\sigma^2}\Ktil = I$.

\paragraph{The obstacle.} Both $T_1$ and $T_2$ contain $\Ktil^{-1}$, an $N\times
N$ dense inverse. A direct solve is $\bigo(N^3)$; the trace is worse. To make
these tractable we need structure in $K$.

\section{Step 2. The EFGP factorization}

Fix an equispaced frequency grid $\{\xi_k\}_{k=1}^M\subset\mathbb{R}^d$ with
spacing $h$, chosen to resolve the kernel spectrum to tolerance $\varepsilon$.
Define
\begin{equation}
  F\in\mathbb{C}^{N\times M},\quad F_{j,k} = e^{i2\pi x_j\cdot\xi_k},\qquad
  D = \operatorname{diag}\!\bigl(\sqrt{h^d\,\hat k_\theta(\xi_k)}\bigr),\qquad
  \Phi := FD.
\end{equation}
Quadrature on the kernel Bochner integral gives the two approximations we need:
\begin{equation}
  K_\theta \;\approx\; \Phi\Phi^* \;=\; FD^2F^*,\qquad
  \partial_\theta K_\theta \;\approx\; F\,D'_\theta\,F^*,\qquad
  D'_\theta := \operatorname{diag}\!\bigl(h^d\,\partial_\theta\hat k_\theta(\xi_k)\bigr).
  \label{eq:efgp}
\end{equation}
Both approximations share the \emph{same} $F$; only the spectral diagonal
changes. This is the structural fact that will let us reduce every $N$-space
expression to an $M$-space expression.

Two auxiliary $M\times M$ operators appear repeatedly:
\begin{equation}
  T \;:=\; F^*F\quad(\text{multi-level Toeplitz, $\bigo(M\log M)$ matvec}),\qquad
  \Amean \;:=\; \sigma^2 I + DTD \;=\; \sigma^2 I + \Phi^*\Phi.
  \label{eq:AT}
\end{equation}
We will apply CG only against $\Amean$, never against the $N\times N$ $\Ktil$.

\section{Step 3. Solving the mean: Woodbury moves the inverse to $M$-space}

Because $\Ktil = \sigma^2 I + \Phi\Phi^*$ is a rank-$M$ update of a diagonal, the
Sherman--Morrison--Woodbury identity applies directly:
\begin{equation}
  \Ktil^{-1}
  \;=\; \sigma^{-2}\!\left(I - \Phi\,(\sigma^2 I + \Phi^*\Phi)^{-1}\,\Phi^*\right)
  \;=\; \sigma^{-2}\!\left(I - \Phi\,\Amean^{-1}\,\Phi^*\right).
  \label{eq:woodbury}
\end{equation}
Applying both sides to $y$,
\begin{equation}
  \alpha \;=\; \Ktil^{-1}y \;=\; \sigma^{-2}\bigl(y - \Phi\beta_0\bigr)
                       \;=\; \sigma^{-2}\bigl(y - FD\beta_0\bigr),\qquad
  \beta_0 \;:=\; \Amean^{-1}\,DF^*y.
  \label{eq:alpha_beta0}
\end{equation}
This reduces the mean to one NUFFT ($z := F^*y$) followed by one $M$-space CG
solve ($\Amean\beta_0 = Dz$).

\paragraph{A useful reorganization.} The equation defining $\beta_0$,
$\Amean\beta_0 = DF^*y$, expands (using $\Amean = \sigma^2 I + DTD$) as
$\sigma^2\beta_0 + DTD\beta_0 = DF^*y$. Rearrange:
\begin{equation}
  D\bigl(F^*y - T\,D\beta_0\bigr) \;=\; \sigma^2\beta_0.
  \label{eq:mean_rearrange}
\end{equation}
Separately, applying $F^*$ to \eqref{eq:alpha_beta0} and using $T = F^*F$
(eq.~\eqref{eq:AT}):
\begin{equation}
  F^*\alpha \;=\; \sigma^{-2}\bigl(F^*y - T\,D\beta_0\bigr).
  \label{eq:Fadj_alpha}
\end{equation}
Multiplying \eqref{eq:Fadj_alpha} by $D$ and comparing with
\eqref{eq:mean_rearrange}: $\sigma^2\,DF^*\alpha = D(F^*y - TD\beta_0) =
\sigma^2\beta_0$, so
\begin{equation}
  \boxed{\;D\,F^*\alpha \;=\; \beta_0\;}\qquad\text{(exactly, at CG convergence)}.
  \label{eq:key_mean}
\end{equation}
Equivalently, $F^*\alpha = \beta_0/D$ entrywise on $\operatorname{supp}(D)$.
This identity is the reason the $\beta$-form exists: the quantity $F^*\alpha$,
which would otherwise require a $\sigma^{-2}$-amplified combination of $F^*y$
and $T(D\beta_0)$, is already sitting in memory as $\beta_0$ --- up to a
pointwise division by $D$.

\section{Step 4. $T_2$ in $M$-space}

\subsection*{Kernel hypers}

Substitute \eqref{eq:efgp} into \eqref{eq:T2alpha}:
\begin{equation}
  T_2(\theta) \;=\; \alpha^*\bigl(F\,D'_\theta\,F^*\bigr)\alpha
              \;=\; \langle F^*\alpha,\; D'_\theta\,F^*\alpha\rangle.
  \label{eq:T2_Fform}
\end{equation}
We now have two algebraically equivalent routes.

\paragraph{Route A (naive Woodbury / residual form).} Use
eq.~\eqref{eq:Fadj_alpha} literally, $F^*\alpha = \sigma^{-2}(F^*y - TD\beta_0)$.
Every evaluation carries a $\sigma^{-2}$ prefactor. Under an inexact CG solve
with residual $r := \Amean\beta_0 - DF^*y$, the quantity evaluated is
$F^*\alpha + r/(\sigma^2 D)$ (on the $D$-support). Because
\eqref{eq:T2_Fform} is \emph{quadratic} in $F^*\alpha$, the error contains a
$\|r\|^2/\sigma^4$ term that can flip the sign of the gradient when $\sigma^2$
is small.

\paragraph{Route B (the $\beta$-form).} Use the identity \eqref{eq:key_mean}:
$F^*\alpha = \beta_0/D$. Then \eqref{eq:T2_Fform} becomes
\begin{equation}
  \boxed{\;T_2(\theta) \;=\; \bigl\langle \beta_0,\; (D'_\theta/|D|^2)\,\beta_0\bigr\rangle.\;}
  \label{eq:T2beta}
\end{equation}
This uses only the CG output $\beta_0$. Under inexact CG the error in $\beta_0$
enters linearly (through a single factor of $\beta_0$ on each side of the inner
product), with no $1/\sigma$ prefactor. The $1/|D|^2$ weight is bounded on the
EFGP grid by construction (the grid is truncated where $\hat k(\xi)$ falls
below $\varepsilon$).

\subsection*{Noise}

Since $\partial_{\sigma^2}\Ktil = I$, we have $T_2(\sigma^2) = \|\alpha\|^2$.
Expanding the defining relation $\sigma^2\alpha = y - FD\beta_0$ directly,
\begin{equation}
  \boxed{\;\|\alpha\|^2 \;=\; \sigma^{-4}\bigl(\|y\|^2 - 2\,\Re\langle\beta_0,\,DF^*y\rangle + \langle\beta_0,\,DTD\,\beta_0\rangle\bigr).\;}
  \label{eq:T2noise}
\end{equation}
Crucially this expansion does \emph{not} invoke
$\Amean\beta_0 = DF^*y$, so it is CG-robust: all three scalars are $M$-dots on
vectors already computed during the mean solve ($F^*y$, $D\beta_0$,
$T(D\beta_0)$).

\subsection*{Variance $\sigma_f^2$ (closed form)}

For SE/Mat\'ern kernels, $\hat k$ is linear in $\sigma_f^2$, so
$D^2 = \sigma_f^2\cdot(\text{shape})$ and $D'_{\sigma_f^2} = D^2/\sigma_f^2$.
Substituting into \eqref{eq:T2beta}:
\[
  T_2(\sigma_f^2) \;=\; \|\beta_0\|^2/\sigma_f^2.
\]

\section{Step 5. $T_1$: trace in $M$-space}

\subsection*{Kernel hypers}

By cyclicity of trace and \eqref{eq:efgp},
\begin{equation}
  T_1(\theta) \;=\; \tr\!\bigl(\Ktil^{-1}\,F D'_\theta F^*\bigr)
              \;=\; \tr\!\bigl(F^*\Ktil^{-1}F\,\cdot\,D'_\theta\bigr)
              \;=\; \tr\!\bigl(B\,D'_\theta\bigr),\qquad
  B \;:=\; F^*\,\Ktil^{-1}\,F \;\in\;\mathbb{C}^{M\times M}.
  \label{eq:Bdef}
\end{equation}
The trace has moved into $M$-space. Now apply Woodbury \eqref{eq:woodbury} to
expand $B$:
\begin{equation}
  B \;=\; \sigma^{-2}\bigl(F^*F - F^*\Phi\,\Amean^{-1}\,\Phi^*F\bigr)
     \;=\; \sigma^{-2}\bigl(T - T D\,\Amean^{-1}\,D T\bigr).
  \label{eq:Bform}
\end{equation}

\paragraph{Hutchinson.} For a Rademacher probe $v\in\{\pm1\}^M$,
$\mathbb{E}_v[v^*\,BD'_\theta\,v] = \tr(BD'_\theta) = T_1(\theta)$. Let
\[
  u^\theta \;:=\; D'_\theta\,v,\qquad
  X^\theta \;:=\; \Amean^{-1}\bigl(D T u^\theta\bigr).
\]
Then \eqref{eq:Bform} gives
\begin{equation}
  B\,D'_\theta\,v \;=\; \sigma^{-2}\bigl(T u^\theta - T D\,X^\theta\bigr).
  \label{eq:BDv}
\end{equation}

\paragraph{Route A (naive Woodbury).} Evaluate \eqref{eq:BDv} literally and
inner-product with $v$. The $\sigma^{-2}$ prefactor amplifies the CG residual of
$X^\theta$.

\paragraph{Route B (the $\beta$-form).} Apply the same reorganization as in
Step~3 to the defining equation $\Amean X^\theta = DTu^\theta$:
\begin{equation}
  \sigma^2 X^\theta + DTD\,X^\theta \;=\; DT u^\theta
  \quad\Longrightarrow\quad
  D\bigl(T u^\theta - T D\,X^\theta\bigr) \;=\; \sigma^2\,X^\theta.
  \label{eq:trace_rearrange}
\end{equation}
Multiplying \eqref{eq:BDv} by $\sigma^2 D$ and comparing with
\eqref{eq:trace_rearrange}: $\sigma^2\,D\,BD'_\theta v = D(Tu^\theta -
TDX^\theta) = \sigma^2 X^\theta$, so
\begin{equation}
  \boxed{\;B\,D'_\theta\,v \;=\; X^\theta/D\;}\qquad\text{(exactly, at CG convergence)}.
  \label{eq:key_trace}
\end{equation}
Substituting \eqref{eq:key_trace} into the Hutchinson estimator,
\begin{equation}
  T_1(\theta) \;\approx\; \frac{1}{J}\sum_{j=1}^{J}\bigl\langle v_j,\; B D'_\theta v_j\bigr\rangle
              \;=\; \frac{1}{J}\sum_{j=1}^{J}\bigl\langle v_j,\; X^\theta_j/D\bigr\rangle.
  \label{eq:T1beta}
\end{equation}
One batched CG solve against $\Amean$ produces the $\{X^\theta_j\}$; the rest
is a Hadamard divide and an inner product. No NUFFT and no Toeplitz application
enter \eqref{eq:T1beta} after the right-hand side is formed.

\paragraph{Why this is lower-variance than a naive $N$-space Hutchinson.}
A data-space Hutchinson estimator $z^*\Ktil^{-1}(\partial_\theta K)z$ with
$z\in\{\pm1\}^N$ has its Hutchinson variance proportional to the squared
Frobenius norm of the off-diagonal of $\Ktil^{-1}\partial_\theta K$. The same
trace in $M$-space is $v^*\,BD'_\theta\,v$ with $v\in\{\pm1\}^M$, and
$BD'_\theta$ is $M\times M$. Since $\partial_\theta K \approx FD'_\theta F^*$
has rank at most $M$ by construction, the $N\!-\!M$ directions where the
$N$-space matrix acts trivially still contribute Hutchinson variance. The
$M$-space probe sees only the informative subspace: single-probe variance drops
by roughly the rank ratio $M/N$.

\subsection*{Noise}

From Woodbury \eqref{eq:woodbury} applied to the trace,
\begin{equation}
  T_1(\sigma^2) \;=\; \tr(\Ktil^{-1})
           \;=\; \sigma^{-2}\bigl(N - \tr(\Amean^{-1}\,DTD)\bigr),
  \label{eq:T1noise}
\end{equation}
with the $M$-space trace estimated by Hutchinson on $\Amean^{-1}DTD$ using
probes $V_n\in\{\pm1\}^{M\times J}$ solved in the same batched CG.

\subsection*{Variance $\sigma_f^2$ (closed form)}

Since $K_\theta = \sigma_f^2\cdot\partial_{\sigma_f^2}K_\theta$,
\[
  T_1(\sigma_f^2) \;=\; \frac{1}{\sigma_f^2}\,\tr(\Ktil^{-1}K)
              \;=\; \frac{N - \sigma^2\,T_1(\sigma^2)}{\sigma_f^2}.
\]
No separate Hutchinson call.

\section{Step 6. Why the $\beta$-form everywhere}

Both $T_2$ and $T_1$ land on the same structural choice: to express the needed
quantity ($F^*\alpha$ or $BD'_\theta v$), either
\begin{itemize}
  \item[(A)] evaluate the Woodbury expression literally, which carries a
    $\sigma^{-2}$ prefactor on the CG residual; or
  \item[(B)] apply the reorganization $D\cdot(\text{Woodbury expression}) =
    \sigma^2\cdot(\text{CG output})$ and divide by $D$.
\end{itemize}
The two are algebraically identical at CG convergence. Under inexact CG with
residual $\delta$, Route~A contributes $\|\delta\|/\sigma^2$; in $T_2$, where
this sits on \emph{both sides} of a quadratic form, it becomes
$\|\delta\|^2/\sigma^4$ and can flip the gradient sign at small $\sigma^2$.
Route~B replaces the $1/\sigma^2$ amplification with a $1/D$ division that is
bounded on the EFGP grid (the grid is truncated where
$\hat k(\xi) < \varepsilon$, so $|D|^2 \gtrsim \varepsilon$, and the amplification
is $\bigo(1/\varepsilon)$ rather than $\bigo(1/\sigma^2)$).
This is the entire motivation for the $\beta$-form.

\section{Step 7. Algorithm}

Let $\theta_1,\dots,\theta_{K-1}$ be the kernel hypers excluding $\sigma_f^2$.
Precompute $F$, $D$, $\{D'_{\theta_i}\}$, and a Jacobi preconditioner for
$\Amean$.

\begin{algorithm}[h]
\caption{EFGP log-marginal gradient: one NUFFT, two CG calls.}
\begin{algorithmic}[1]
\State $z \leftarrow F^*y$ \Comment{\textbf{only NUFFT in the whole gradient}}
\State Solve $\Amean\,\beta_0 = D z$ by CG. \hfill\textbf{(CG \#1, mean)}
\State $w \leftarrow T(D\beta_0)$ \Comment{Toeplitz matvec, reused}
\Statex\hrulefill
\State \emph{Term $T_2$ (all $M$-space, Route B where relevant):}
\State \quad kernel hypers: $T_2(\theta_i) \leftarrow \langle\beta_0,\,(D'_{\theta_i}/|D|^2)\,\beta_0\rangle$  \hfill [eq.~\eqref{eq:T2beta}]
\State \quad variance:      $T_2(\sigma_f^2) \leftarrow \|\beta_0\|^2/\sigma_f^2$
\State \quad noise:         $T_2(\sigma^2) \leftarrow \sigma^{-4}\bigl(\|y\|^2 - 2\,\Re\langle\beta_0, Dz\rangle + \langle\beta_0, Dw\rangle\bigr)$ \hfill [eq.~\eqref{eq:T2noise}]
\Statex\hrulefill
\State \emph{Term $T_1$:} draw $V\in\{\pm1\}^{M\times J}$, $V_n\in\{\pm1\}^{M\times J}$.
\State \quad kernel: for each $\theta_i$, $\mathrm{rhs}^i \leftarrow D\,T\,(D'_{\theta_i}V)$
\State \quad noise:  $\mathrm{rhs}^n \leftarrow D\,T\,(D V_n)$
\State \quad stack RHSs; solve $\Amean X = \mathrm{rhs}$ by batched CG. \hfill\textbf{(CG \#2, trace)}
\State \quad kernel ($\beta$-form): $T_1(\theta_i) \leftarrow \mathrm{mean}_j\,\langle V_j,\,X^i_j/D\rangle$  \hfill [eq.~\eqref{eq:T1beta}]
\State \quad noise:  $T_1(\sigma^2) \leftarrow \sigma^{-2}\bigl(N - \mathrm{mean}_j\,\langle (V_n)_j,\,X^n_j\rangle\bigr)$  \hfill [eq.~\eqref{eq:T1noise}]
\State \quad variance: $T_1(\sigma_f^2) \leftarrow (N - \sigma^2\,T_1(\sigma^2))/\sigma_f^2$
\Statex\hrulefill
\State \textbf{Return} $\nabla_\theta\log p = \tfrac12(T_2 - T_1)$ for each hyper.
\end{algorithmic}
\end{algorithm}

\paragraph{Complexity.} One NUFFT ($\bigo(N + M\log M)$). Two CG calls against
$\Amean$, each $\bigo(M\log M)$ per matvec: the mean with one RHS, the trace
with $(K\!-\!1)\!+\!1 = K$ blocks of $J$ RHSs stacked. All $T_2$ work is
$\bigo(KM)$. No NUFFT and no $N$-space linear algebra appear inside the trace
loop.

\end{document}
